\section{Principle II: Force}

To the principle of Truth, Mr. Maurras adds the principle of Force, without which the former is impotent. Yet, Force must be guided by sound judgment. 

There are despotic weaknesses, malicious stupidities, and losers worthy of being so, as there are beneficent winners, heroes of energy and power to whom humanity owes immense progress, giants of health and strength who have merited the blessing of the past and the future. Force in itself, bared of its adventitious and incidental characteristics, force which is still in the service of neither good nor evil, naked force and by itself a good, and very precious and very great since it is the expression of the activity of being. It is idiotic to want to ignore its benefits.

In order to deter harmful sophism, one cannot grow tired of repeating it: force by itself, reduced to itself, is a good. That does not mean that it always does good or there is not a greater good.

Good in itself, it is also capable of the greatest benefits like defending the country, punishing crime, avenging honor, or protecting innocence.

But as it is capable of everything, it needs a rule as the primary guaranty and, put into the service of the best cause, an order. Order contributes to render it entirely and completely effective. But order also restrains it in the service of what it claims to serve; order prevents it from turning against what is dear to it in spite of itself. All disordered force is exposed to this peril.

It is up to Reason to moderate force, that is to say, the sense of measure and intellectual proportions, there is this civic sense which adds to the higher orders of the spirit I don't know what principle of cordiality, bonhomie, I would dare to say, of charity which does not desire the death of the guilty, but which desires, and very much desires, his rectification.

There remains also that clarity, that sincerity, the natural continuation of the rectitude of the spirit which permits neither hypocrisy nor disguise, but which marches straight ahead of ifself, head high and chest bare, this tranquil serenity which belongs to those who have freely assumed noble missions.

The theories of force are not at all in contradiction with the doctrine of solidarity and, anyway, human mutual aid needs to be strong to protect itself or to be protected against violence.

If there may be in the world force which is of value, it appears indispensable for it to be strong; if there is something else, if, as we think, there is something better, and much better, it is still more necessary to be strong and powerful to save or to develop these true goods.



\flrightit{Posted on 2012-10-07 by Cologero }

\begin{center}* * *\end{center}

\begin{footnotesize}\begin{sffamily}



\texttt{Audrey on 2012-10-07 at 23:02 said: }

I do find your posts to be a good mirror of self-reflection, contributing to the `educatedness' of my self-observations and regulations of behavior. I would say, that real education should approach, at least in pary, a similar mode, as opposed to stuffing slots in people's heads with ideas and fighting fr intellectual real estate. Namely, paragraph \#4 from the bottom, I have always wondered why every time I am asked about punishment under UCMJ, such as rank degradation, or jail, I always request instead rectificator measures. However, in the Army, this is considered a weakness. I guess their idea of teamwork, is more about pulling like a donkey, and not about actual collaboration.


\hfill

\texttt{Logres on 2012-10-08 at 05:22 said: }

Bishop Joubert (as quoted by M. Arnold): ``Force \& Right rule the world – Force till Right is ready". Maurras' style and manner of thought comes across very strongly, even in translation. Thank you.


\hfill

\texttt{Sati Shankar on 2012-10-11 at 22:48 said: }

Force,is in fact a double edge sword,neutral it itself as such, in the hands of some one capable in using it.Now it depends on the its user that it is used to protect or induce Good. History has enough to cite its misuses. The question is, whether to engineer (by force) a society where corrections will not be needed which depends, dangerously upon the user of force, Or to let the Tradition work so that no engineering is needed at all. Force kills spontaneity and freedom, most often.

The second point relates to the Force as inner energy, that is indispensable for an individual to attain his goal.


\hfill

\texttt{Graham on 2012-10-14 at 01:21 said: }

``… neutral in itself …"

Maurras says that force is good in itself. That makes sense, since ``it is the expression of the activity of being". This closely parallels Evola's remarks about virtue and power which are discussed elsewhere on Gornahoor.


\end{sffamily}\end{footnotesize}
